\section{Scope of the curvature condition}\label{sec:structure}
\subsection{Uniform smoothness and the gradient-norm model}
Since $f$ is $C^2$ near a compact $Q$, the uniform bound
$L=\max_{x\in Q}\|\nabla^2f(x)\|$ is finite. Our assumption is useful
through its parameters and the resulting local model, not because it admits
functions without a finite smoothness constant on this fixed $Q$.
Moreover,
\[
 G(x)\le D\|\nabla f(x)\|,
 \qquad \|\nabla^2f(x)\|\le L_0+(L_1D)\|\nabla f(x)\|.
\]
Thus gap-dependent smoothness implies a norm-dependent bound on $Q$.
The converse with the same zero constant term fails: $f(t)=e^t$ on
$[0,R]$ satisfies $|f''|=|f'|$, whereas $G(0)=0$ and $f''(0)=1$.
Every admissible gap-dependent pair therefore has $L_0\ge1$ in that example.
In general $L_0\ge\|\nabla^2f(x^*)\|$.

\begin{proposition}[A restriction of the pure-gap regime]\label{prop:rigidity}
Suppose \eqref{eq:gap_smoothness} holds with $L_0=0$.
If $\nabla f(z)=0$ at any $z\in Q$, then $\nabla f=0$ on $Q$ and
$f$ is constant on $Q$. Convexity or quasar-convexity of $f$ is not needed.
\end{proposition}
\begin{proof}
Fix $y\in Q$ and set $u(t)=\|\nabla f(z+t(y-z))\|$.
The fundamental theorem of calculus and $G\le D\|\nabla f\|$ give
\[
 u(t)\le L_1D\|y-z\|\int_0^t u(s)\,ds.
\]
Gronwall's inequality yields $u=0$. Varying $y$ proves the assertion.
\end{proof}
Consequently, a nonconstant problem in the pure-gap regime cannot have
a minimizer with zero gradient. In particular, an interior minimizer of a
full-dimensional feasible set is excluded. This restriction should be kept
in mind when interpreting Corollary~\ref{cor:pure_gap_linear}.

\subsection{Connection to existing growth conditions}
For convex $f$, write
$D_f(y,x)=f(y)-f(x)-\langle\nabla f(x),y-x\rangle$.
If $L_0=0<L_1$, Lemma~\ref{lem:transport} implies, for $0\le\alpha\le1$,
\[
 D_f(x+\alpha d,x)
 \le \frac{M\alpha^2}{2}G(x),
 \qquad M=L_1D^2e^{L_1D^2}.
\]
Indeed, bound the exponential in \eqref{eq:upper_model} by
$e^{L_1D^2}$ and integrate $\alpha-t$.
This is the strong $(M,1)$-growth condition used in~\cite{wirth}.
The reduction is conservative: it introduces an exponential dependence on
$L_1D^2$, whereas Corollary~\ref{cor:pure_gap_linear} uses an explicit
step lower bound of order $\min\{1,(L_1D^2)^{-1}\}$.
This comparison concerns the displayed conversion of constants, not a
lower bound on what all other analyses or algorithms can achieve.

\subsection{A multidimensional class with explicit parameters}
Let $a_i\in\mathbb R^n$, $y_i\ge0$, and $B=\max_i\|a_i\|>0$.
Consider the Poisson negative log-likelihood, up to data-only constants,
\begin{equation}\label{eq:poisson}
 f(x)=\frac1m\sum_{i=1}^m\bigl(e^{a_i^Tx}-y_i a_i^Tx\bigr),
 \qquad 0\in Q.
\end{equation}
For integer $y_i$ this is a count-regression objective; the following
bound also holds for nonnegative real responses.

\begin{proposition}[Parameters from the observed data]\label{prop:poisson}
On any compact convex $Q$ containing zero, \eqref{eq:poisson} satisfies
\eqref{eq:gap_smoothness} with
\begin{equation}\label{eq:poisson_constants}
 L_0=B^2\left(1+\frac{1}{2m}\sum_i y_i^2\right),
 \qquad L_1=B^2,\qquad \gamma=1.
\end{equation}
These constants do not require a minimizer or a minimum value.
If $Q=\{x:\|x\|\le R\}$, valid norm-dependent constants are
$K_0=L_0$, $K_1=RL_1$, and a uniform Hessian bound is
\begin{equation}\label{eq:poisson_uniform}
 L=e^{BR}\|A\|_2^2/m,
\end{equation}
where the rows of $A$ are $a_i^T$. These are valid bounds, not claims
of smallest possible constants.
\end{proposition}
\begin{proof}
For $y\ge0$ and $t\in\mathbb R$,
\begin{equation}\label{eq:poisson_scalar}
 e^t\le 1+y^2/2+t(e^t-y).
\end{equation}
For $t\le0$, $(1-t)e^t+yt\le1$.
For $t\ge0$, the derivative of $(1-t)e^t$ is $-te^t\le-t$, so
$(1-t)e^t+yt\le1-t^2/2+yt\le1+y^2/2$.
Also $\|\nabla^2f(x)\|\le B^2m^{-1}\sum_i e^{a_i^Tx}$.
Summing \eqref{eq:poisson_scalar} therefore gives
\[
 \|\nabla^2f(x)\|\le L_0+L_1\langle\nabla f(x),x\rangle
 \le L_0+L_1G(x),
\]
where the last step uses $0\in Q$.
On a ball, $\langle\nabla f(x),x\rangle\le R\|\nabla f(x)\|$.
Finally $a_i^Tx\le BR$ implies
$\nabla^2f(x)\preceq e^{BR}A^TA/m$.
\end{proof}
The ball's linear oracle is $s=-R\nabla f(x)/\|\nabla f(x)\|$
when the gradient is nonzero. It generally differs from the minimizer
of \eqref{eq:poisson}. Thus these problems avoid the endpoint degeneracy
of the scalar examples below. The model and parameter bounds are global
on $Q$; the numerical experiments only sample their consequences.
